\clearpage
\VerdierPhysicalPageAnchor{0031}
\VerdierSemanticAnchor{intro-8-etale-lefschetz}
在 étale 拓扑的框架中，也有类似的 Lefschetz 公式 [12]。
\VerdierSemanticAnchor{intro-8-frobenius-L}
把该公式应用于有限域上光滑完备代数曲线的 Frobenius 自同态，并在曲线上
取一个关于 étale 拓扑可构造的层，便得到了广义 $L$ 函数有理性的证明
[11]。应当指出，该证明所用的 Lefschetz 公式中，系数层确实带有奇点。

\VerdierSemanticAnchor{intro-9}
\noindent\textbf{9.}\quad
在拓扑 Lefschetz 公式的情形中，我们已经看到复形
\VerdierSemanticAnchor{intro-9-perfect-complex}
\[
\mathrm{RH}(X,F)
\]
是完美的，从而可以定义下列自同态的 Lefschetz 数：
\VerdierSemanticAnchor{intro-9-rh-phi}
\[
\mathrm{RH}(\Phi)\colon
\mathrm{RH}(X,F)\longrightarrow\mathrm{RH}(X,F).
\]

\VerdierSemanticAnchor{intro-9-cohomology-objects}
这格外鲜明地说明了为什么不能用一个复形的上同调对象族来取代该复形。实际
上，只有在极为严格的条件下，才能定义自同态
\VerdierSemanticAnchor{intro-9-hi-phi}
\[
\mathrm{H}^i(\Phi)\colon
\mathrm{H}^i(X,F)\longrightarrow\mathrm{H}^i(X,F)
\]
的迹。

\VerdierSemanticAnchor{intro-9-regular-trace}
例如，当环 $A$ 的全局上同调维数有限，也就是 $A$ 正则时，可通过取有限
投射分解来定义 $\mathrm{H}^i(\Phi)$ 的迹。此时有
\VerdierSemanticAnchor{intro-9-trace-formula}
\[
\chi\bigl(\mathrm{RH}(\Phi)\bigr)
=\sum_i(-1)^i\operatorname{Tr}\bigl(\mathrm{H}^i(\Phi)\bigr).
\]

\VerdierSemanticAnchor{intro-9-nonregular-rings}
然而，如前所见，尤其在 étale 拓扑的层论中，必须处理
$\mathbb{Z}/p^\ell\mathbb{Z}$（$\ell\geq 2$）这类非正则环。

\VerdierSemanticAnchor{intro-9-k-theory}
完美复形这一概念在代数几何中同样有用：K 理论与 Riemann--Roch 定理 [14]。

\VerdierSemanticAnchor{intro-10}
\noindent\textbf{10.}\quad
综上，目前看来，导出范畴的语言是表述并证明上述理论中基本结果所不可缺少
的工具。

\VerdierSemanticAnchor{intro-10-additive-scope}
本文限于加性同调代数：阿贝尔范畴与加性函子。一旦涉及 Riemann--Roch
定理 [14]，这一框架便不再充分。
\VerdierSemanticAnchor{intro-10-omissions}
我们也没有研究可由三角范畴构造的各种类群或 Grothendieck 群。最后，相对
同调的问题也被留在本文范围之外，尽管它们同样可以在三角范畴框架中处理
[15]。
